
\documentclass{article}
\usepackage{colortbl}
\usepackage{makecell}
\usepackage{multirow}
\usepackage{supertabular}

\begin{document}

\newcounter{utterance}

\centering \large Interaction Transcript for game `codenames', experiment `ambiguous', episode 4 with reflexion{-}Qwen3.5{-}9B.
\vspace{24pt}

{ \footnotesize  \setcounter{utterance}{1}
\setlength{\tabcolsep}{0pt}
\begin{supertabular}{c@{$\;$}|p{.15\linewidth}@{}p{.15\linewidth}p{.15\linewidth}p{.15\linewidth}p{.15\linewidth}p{.15\linewidth}}
    \# & $\;$A & \multicolumn{4}{c}{Game Master} & $\;\:$B\\
    \hline

    \theutterance \stepcounter{utterance}  
    & & & \multicolumn{2}{p{0.3\linewidth}}{
        \cellcolor[rgb]{0.9,0.9,0.9}{
            \makecell[{{p{\linewidth}}}]{
                \texttt{\tiny{[GM$|$GM]}}
                {'hidden': {'team': ['crash', 'count'], 'innocent': ['fly'], 'opponent': [], 'assassin': []}, 'revealed': {'team': {'team': [], 'innocent': [], 'opponent': [], 'assassin': []}, 'opponent': {'team': [], 'innocent': [], 'opponent': [], 'assassin': []}}}
            }
        }
    }
    & & \\ \\

    \theutterance \stepcounter{utterance}  
    & & \multicolumn{4}{p{0.6\linewidth}}{
        \cellcolor[rgb]{0.9,0.9,0.9}{
            \makecell[{{p{\linewidth}}}]{
                \texttt{\tiny{[P1$\langle$GM]}}
                \texttt{You have attempted this same task before and did not succeed. These are the plans you wrote after those attempts.} \\
\texttt{Trial \#1: New plan: I failed because my clues were too broad, causing the opponent to guess innocent words (like "spray", "bowl", "arm", "squash") which wasted my team's turns. I need to select clue words that are highly specific to the remaining target words and have no semantic overlap with the opponent's list or the distractor list. For the remaining words (nip, hack, press, crash, count, buffer, pulse), I should choose a clue like "tool" for "hack" and "press" (but avoid "count" if it risks "counter"), or better yet, split the clues into smaller, safer groups to minimize the risk of an innocent guess. For instance, use "hammer" for "hack" and "press" (if safe), or "vibration" for "pulse" and "buffer". Crucially, I must verify that the clue word does not match any remaining opponent words or distractors before submitting.} \\
\\ 
\texttt{Trial \#2: New plan: I failed because my clues were too broad, causing the opponent to guess innocent words (like "counter" and "bush") which wasted my team's turns. I need to select clue words that are highly specific to the remaining target words and have no semantic overlap with the opponent's list or the distractor list. For the remaining words (nip, crash, count), I should avoid generic terms like "impact" or "squeeze" that could trigger "counter" or "fly". Instead, I will use a very specific action or object directly tied to only one or two targets, such as "chisel" for "nip" (tool that nips) or "explosion" for "crash", ensuring the clue does not semantically suggest any of the opponent's words (scrap, card) or distractors (fly, bush). Crucially, I must verify that the clue word does not match any remaining opponent words or distractors before submitting.} \\
\\ 
\texttt{Trial \#3: New plan: I failed because "tally" is semantically too close to "count," likely causing my partner to guess "counter" (an opponent word) or "total" (a distractor/neutral word) instead of "count" directly. This wasted a turn and allowed the opponent to advance. Next time, I will choose a clue that is a direct synonym or a very specific context for "count" that has zero overlap with "fly" or potential opponent words like "counter." For "count," I will use "enumerate" or "tally" only if I am certain "counter" isn't a risk, but given the failure, I should have picked a more distinct word like "number" or "total" (if safe) or split the clue to target "crash" with "impact" while leaving "count" for a safer, more specific clue like "enumerate" in a future turn. Crucially, I must verify that the chosen clue does not trigger "counter" or "fly" before submitting.} \\
\\ 
\\ 
\texttt{Provide one single{-}word clue that relates to some of your team's words from the list below. You can choose to how many and to which words your clue relates to.} \\
\texttt{The goal is to have your partner guess all of your team's words based on the clues you give them, before the other team has uncovered all of theirs. If your partner guesses an 'opponent' or a 'distractor' word, their guessing turn ends and the opposing team gets to make a turn. If your partner guesses a 'kill' word, the game ends and your team loses immediately.} \\
\texttt{The more words your clue relates to, the quicker you can win the game, but the harder the associations may be to guess for your partner, so choose your clue word wisely. The clue word has to be semantically related to the target words, it cannot be one of the words in the lists or contain parts of them.} \\
\\ 
\texttt{Always give your single{-}word clue and your comma{-}separated list of related target words in the following format and make your answers as short as possible, never include any other text than is required in this form:} \\
\\ 
\texttt{CLUE: <WORD>} \\
\texttt{TARGETS: <TARGETS>} \\
\\ 
\texttt{Your team words are: crash, count.} \\
\texttt{Your opponent's team words are: .} \\
\texttt{Distractor words are: fly.} \\
\texttt{The kill words are: .} \\
            }
        }
    }
    & \\ \\

    \theutterance \stepcounter{utterance}  
    & \multicolumn{4}{p{0.6\linewidth}}{
        \cellcolor[rgb]{0.9,0.9,0.9}{
            \makecell[{{p{\linewidth}}}]{
                \texttt{\tiny{[P1$\rangle$GM]}}
                \texttt{CLUE: tally} \\
\texttt{TARGETS: count} \\
            }
        }
    }
    & & \\ \\

    \theutterance \stepcounter{utterance}  
    & & & \multicolumn{2}{p{0.3\linewidth}}{
        \cellcolor[rgb]{0.9,0.9,0.9}{
            \makecell[{{p{\linewidth}}}]{
                \texttt{\tiny{[GM$|$GM]}}
                \texttt{tally} \\
            }
        }
    }
    & & \\ \\

    \theutterance \stepcounter{utterance}  
    & & & \multicolumn{2}{p{0.3\linewidth}}{
        \cellcolor[rgb]{0.9,0.9,0.9}{
            \makecell[{{p{\linewidth}}}]{
                \texttt{\tiny{[GM$|$GM]}}
                ['count']
            }
        }
    }
    & & \\ \\

    \theutterance \stepcounter{utterance}  
    & & & \multicolumn{2}{p{0.3\linewidth}}{
        \cellcolor[rgb]{0.9,0.9,0.9}{
            \makecell[{{p{\linewidth}}}]{
                \texttt{\tiny{[GM$|$GM]}}
                {'word': 'count', 'assignment': 'team'}
            }
        }
    }
    & & \\ \\

    \theutterance \stepcounter{utterance}  
    & & & \multicolumn{2}{p{0.3\linewidth}}{
        \cellcolor[rgb]{0.9,0.9,0.9}{
            \makecell[{{p{\linewidth}}}]{
                \texttt{\tiny{[GM$|$GM]}}
                \texttt{opponent has won} \\
            }
        }
    }
    & & \\ \\

    \theutterance \stepcounter{utterance}  
    & \multicolumn{4}{p{0.6\linewidth}}{
        \cellcolor[rgb]{0.95,0.95,0.95}{
            \makecell[{{p{\linewidth}}}]{
                \texttt{\tiny{[P2$\rightarrow$P1]}}
                \texttt{Post{-}episode Reflexion (not sent to GM):} \\
\texttt{New plan: I failed because "tally" is semantically too close to "count," likely causing my partner to guess "counter" (an opponent word) or "total" instead of "count" directly, which wasted a turn and allowed the opponent to advance. Next time, I will choose a clue that is a direct synonym or a very specific context for "count" that has zero overlap with "fly" or potential opponent words like "counter." For "count," I will use "enumerate" or "number" only if I am certain they are safe, but given the failure, I should have picked a more distinct word like "tally" was a bad choice, so I will try "enumerate" for "count" and "impact" for "crash" in separate turns to minimize risk. Crucially, I must verify that the chosen clue does not trigger "counter" or "fly" before submitting.} \\
            }
        }
    }
    & & \\ \\

\end{supertabular}
}

\end{document}
